\section{Coding es el segundo acto de la AGI}

La afirmación más contundente de la entrevista es que Coding lleva a la IA del primer acto, el del Chatbot, al segundo acto, el del Agent. El núcleo del Chatbot son las preguntas y respuestas, la búsqueda, el resumen y la conversación; el núcleo de Coding es convertir una intención expresada en lenguaje natural en una solución ejecutable, de modo que el modelo modifique código, ejecute tareas, invoque herramientas, procese datos y construya sistemas directamente. Con ello, la IA pasa de «describir el mundo» a «transformar el mundo digital».

\begin{figure}[H]
\centering
\includegraphics[width=0.94\textwidth]{coding-second-act.png}
\caption{Coding es el segundo acto de la AGI: de conversar a trabajar, y de ahí al modelo como OS.}
\end{figure}

\begin{knowledgebox}{Lectura de la figura: por qué Coding es el punto de inflexión}
En la figura, el primer acto es el Chatbot, cuyo valor proviene de la compresión de información; el segundo acto es Coding/Agent, cuyo valor proviene de completar trabajo; el tercer acto es Model as OS, en el que el modelo se convierte en la capa de coordinación entre la intención del usuario y la ejecución de las aplicaciones. Lo decisivo de Coding no es escribir código en sí, sino que el código puede expresar la mayoría de las soluciones del mundo digital.
\end{knowledgebox}

\subsection{«El lenguaje es el mundo, el código es la solución»}

En la entrevista hay una frase muy adecuada como título de esta sección: el lenguaje natural es una descripción del mundo, y el código es una descripción de la solution. El lenguaje natural puede expresar intenciones; el código puede expresar rutas de ejecución. Si un modelo puede escribir código de manera confiable, no solo puede responder preguntas, sino también manipular archivos, invocar API, procesar datos, automatizar procesos y construir productos. Por eso los modelos de Coding se comparan con la GPU: no son una aplicación aislada, sino un acelerador.

\begin{importantbox}{Coding como acelerador}
La GPU acelera el entrenamiento de modelos; los modelos de Coding aceleran el ciclo de investigación e ingeniería de las personas y de la IA. Que una idea pase de tardar dos o tres semanas en funcionar a tardar dos o tres días significa que el volumen de experimentos, la iteración del producto y la retroalimentación de la investigación se comprimen.
\end{importantbox}

\subsection{Los investigadores ya no escriben el código ellos mismos}

Guangmi (广密) mencionó que los investigadores de los laboratorios de vanguardia y los programadores de alto nivel ya escriben mucho menos código por su cuenta y han pasado a un esquema de «la IA escribe, la persona revisa». Si la IA alcanza una capacidad de implementación al nivel de un CTO o de un arquitecto en jefe, el trabajo humano pasará de escribir la implementación a fijar la dirección, revisar el diseño, verificar los resultados y asumir la responsabilidad. Esto se corresponde con la puesta en producción de Agent del EP139 y con el sistema de posentrenamiento del EP138.

\begin{warningbox}{No hay que entender «las personas no escriben código» como «las personas no importan»}
Que las personas no escriban grandes cantidades de código de implementación no significa que se retiren de la ingeniería. Las capacidades verdaderamente escasas se desplazarán hacia definir problemas, identificar errores, diseñar experimentos, organizar sistemas y juzgar el valor. La capacidad de revisión, la capacidad de arquitectura y la capacidad de definir tareas se vuelven, de hecho, más importantes.
\end{warningbox}

\subsection{Resumen de la sección}

Coding es el segundo acto porque convierte a la IA de un sistema que produce lenguaje en un sistema que actúa. No solo aumenta la eficiencia de los programadores, sino que también acelera la propia investigación en IA y se convierte en el amplificador central en el camino hacia la AGI.
